\section{An entropy upper bound for contacts}\label{sec:entropy}

We continue under the assumption that both coordinate signs have positive
probability.  The opposed arrangement and its common blocks are those of
Section~\ref{sec:arrangement}.  We will prove
\begin{proposition}\label{prop:entropy-upper}
Suppose that both coordinate signs have positive directional-transience
probability. In the opposed arrangement of Section~\ref{sec:arrangement},
let $m(J)$ be the expected number of dyadic common-block index intervals
$(2^j,2^{j+1}]$, $1\le j\le J$, containing a departure-site contact,
as in~\eqref{eq:contact-count}. Then
\[
 m(J)=O\!\left(\frac{J}{\log J}\right)\qquad(J\longrightarrow\infty).
\]
\end{proposition}
The argument compares a piece of the opposed arrangement with independent
bridges of prescribed heights.  Contacts are rare under the bridge law.
The cost of this comparison is controlled by the entropy gained when more
independent common-block displacements are added.

\subsection{A randomized comparison with independent bridges}

Fix an integer $n\ge1$.  We want to cover the common blocks with
indices in $(4n,8n]$ by a middle portion of the arrangement.  We will
surround it by two outer portions, each of height at least $n$, so
that any middle contact occurs well inside a finite slab.  A buffer
above all three portions supplies an independent lateral gap.

We randomize the number of blocks in each portion as well as in the
buffer.  Choosing a chunk count uniformly from $n$ possibilities
will offset the cost of specifying its random height.  Randomizing
the buffer count will let us compare its entropy with that of the whole
experiment without paying an additional $\log n$.
Independently of the common-block tape, choose four independent
integer counts
\begin{equation}\label{eq:entropy-count-windows}
 \begin{split}
 U,I_1,I_3&\ \hbox{uniform on }\{n,\ldots,2n-1\},\\
 I_2&\ \hbox{uniform on }\{7n,\ldots,8n-1\}.
 \end{split}
\end{equation}
Read a buffer of $U$ common blocks, followed by three chunks of
$I_1,I_2,I_3$ common blocks.  Let $d_0$ be the lateral gap after the
buffer, namely the sum of its gap increments.  For chunk $i$, let
$A_i$ and $B_i$ be the positive and negative slab lists in their
downward tape order, and let $H_i$ be their common total width.  Write
\[
 H=(H_1,H_2,H_3),\qquad
 \mathcal A=(A_1,A_2,A_3),\qquad
 \mathcal B=(B_1,B_2,B_3).
\]
The three chunks are independent, and $(U,d_0)$ is independent of all
their data.  To see this despite their random boundaries, condition
on the four external counts and use the product law of the common
blocks; the distribution of a chunk depends only on its own count.
The first and third count windows ensure the required outer heights.
The longer middle window ensures that the middle chunk starts before
block $4n$ and ends after block $8n$.

Set $q=d-1$, and let $X_{A,i},X_{B,i}\in\ZZ^q$ be the sums of the
lateral displacements of the words in $A_i,B_i$, measured in the words' actual chronological
directions.  The total positive displacement determines where its
bottom lies relative to its top.  Accordingly, define
\begin{equation}\label{eq:entropy-alignment-variables}
 X_A=\sum_{i=1}^3X_{A,i},\quad X_B=\sum_{i=1}^3X_{B,i},\quad
 Z=d_0+X_A,\quad T=U+I_1+I_2+I_3.
\end{equation}
Here is the placement that explains the alignment variable $Z$ and
will define the same contact event under the actual and reference
laws.  Set $N_*=H_1+H_2+H_3$.  Place the positive chunks
from height $0$ and lateral position $0$, in chronological order
$3,2,1$, reversing the list order within each chunk.  Each word is
still traversed in its original direction.  Place the negative chunks
in chronological order $1,2,3$ from height $N_*-1$ and lateral position
$Z$.  Let $E_{\mathrm{mid}}$ be a departure contact between the two
middle chunks.

For these actual chunk data, this placement reproduces the infinite
opposed arrangement after its buffer, up to translation.  For example,
at the top of the middle chunk the lateral gap, negative minus positive, is
\[
 Z+X_{B,1}-(X_{A,3}+X_{A,2})
   =d_0+X_{A,1}+X_{B,1},
\]
as in the original arrangement; the subsequent list placements also
agree.  From \eqref{eq:entropy-count-windows}, the middle chunk starts
after at most $4n-2$ common blocks and ends after at least $9n$ blocks.
It therefore contains the deterministic common-block interval
$(4n,8n]$.

Denote the actual law of $(H,Z,\mathcal A,\mathcal B)$ by $P_n$.
Define a reference law $Q_n$ as follows.  First sample $(H,Z)$ with its
actual joint marginal.  Given these values, sample the six lists
independently, with respective exact-cut bridge laws
$\mathcal B_{H_i}^+$ and $\mathcal B_{H_i}^-$ from
Lemma~\ref{lem:exact-bridge}.  The reference retains the dependence
between $H$ and $Z$; it changes only the conditional law of the lists.
Apply the same placement rule to its newly sampled lists.  Thus $E_{\mathrm{mid}}$ is one fixed
measurable event in both experiments, and its actual probability
bounds the original contact probability in $(4n,8n]$.

\subsection{The entropy cost of the comparison}

We use entropy and relative entropy in their classical information-theoretic
sense~\cite{Shannon1948,KullbackLeibler1951}, recalling the identities and
bounds needed for the countable path spaces below.

For a countable random variable $V$, write
\[
 \mathcal H(V)=-\sum_v P(V=v)\log P(V=v),
\]
with $0\log0=0$.  Conditional entropy is the average of the corresponding
conditional entropies.  For probability measures $P,Q$ on a countable set,
their relative entropy is
\[
 D(P\Vert Q)=\sum_x P(x)\log\frac{P(x)}{Q(x)},
\]
with value $+\infty$ if $P$ is not absolutely continuous with respect to
$Q$.  Conditional mutual information $I(V;W\mid H)$ is the average,
over $H$, of the relative entropy between the joint conditional law and
the product of its two marginals.

We use the chain rule for relative entropy, obtained by factoring joint
probabilities into marginal and conditional probabilities.  In particular,
if a reference law for $(V,W)$ conditional on $H$ factors between $V$ and
$W$, the averaged conditional relative entropy is at least
$I(V;W\mid H)$.  Nonnegativity of relative entropy follows from Jensen's
inequality applied to $-\log$.  The same factorization, followed by
nonnegativity, shows that applying a measurable map cannot increase
relative entropy.  These facts apply to countable path lists without
assuming that the lists themselves have finite entropy; equivalently one
may first use finite partitions and pass to their increasing limit.

Let $S_1,S_2,\ldots$ be the iid lateral gap increments of the common
blocks, and set
\[
 h_k=\mathcal H(S_1+\cdots+S_k),\qquad k\ge1.
\]
Lemma~\ref{lem:common-blocks} gives $\mathbf E|S_1|<\infty$.  For a
$\ZZ^q$-valued variable with finite first moment, comparison with the
probability weights proportional to $\exp(-|x|/(k+1))$ gives
\begin{equation}\label{eq:gap-entropy-growth}
 h_k\le C+q\log(k+1).
\end{equation}
Indeed the normalizing sum is at most $C_q(k+1)^q$, and the expected
comparison energy is at most $k\mathbf E|S_1|/(k+1)$.  Also $h_k$ is
nondecreasing: condition on the last independent summand, use translation
invariance of entropy, and then remove the conditioning.

By Lemma~\ref{lem:common-blocks}, the sum of the slab radii in a common
block has finite expectation.  Thus each tape's block displacement,
and not only their sum $S$, has a finite first absolute moment.  In
addition the common width has finite mean.  All count, height, and
displacement variables used below therefore have first moments of
order $n$, and their joint and conditional entropies are finite.
No entropy bound for the full path lists is needed.

\begin{lemma}\label{lem:entropy-comparison}
There is a constant $C$, independent of $n$, such that
\[
 D(P_n\Vert Q_n)\le C+h_{14n}-h_n.
\]
\end{lemma}
\begin{proof}
There are two costs to control: imposing the random chunk heights,
and then retaining the aligned starting displacement $Z$.  We first
condition on $H$ alone.  For a specified
pair of slab lists of common total width $h$, let $r$ be the number of
their common positive cumulative widths, including $h$.  The pair is
uniquely parsed into its first $r$ common blocks.  If $r$ belongs to the
count window for chunk $i$, its actual unconditioned mass is
\[
 \frac1n\prod_{a\in A_i}W(a)\prod_{b\in B_i}W(b).
\]
It has zero mass otherwise.  The independent exact-cut bridge pair at
height $h$ has mass given by the same product divided by $u_h^+u_h^-$.
Thus the conditional density of the actual pair, on its support, is
\begin{equation}\label{eq:count-window-density}
 \frac{u_h^+u_h^-}{nP_n(H_i=h)}.
\end{equation}
Common cuts here are counted in the original downward list order,
before any list is reversed for chronological placement.

Let $D_0$ be the averaged conditional relative entropy of
$(\mathcal A,\mathcal B)$ given $H$ with respect to these independent
bridges.  The chunks are independent, so \eqref{eq:count-window-density}
and $u_h^\pm\le1$ give
\[
 D_0\le\sum_{i=1}^3\bigl(\mathcal H(H_i)-\log n\bigr)\le C.
\]
For the last bound, $\mathbf E H_i\le8n\mathbf E K$, and a positive
integer variable with mean $a$ has entropy at most $1+\log(a+1)$,
by comparison with a geometric distribution.  Since the reference
conditional on $H$ factors between the tapes,
\begin{equation}\label{eq:conditional-tape-information}
 I(\mathcal A;\mathcal B\mid H)\le D_0.
\end{equation}

Introducing $Z$ while retaining its actual marginal costs precisely
its conditional mutual information with the lists.  The chain rule
and the independence of the buffer give
\begin{equation}\label{eq:endpoint-information-cost}
 \begin{split}
 D(P_n\Vert Q_n)
 &=D_0+I(Z;\mathcal A,\mathcal B\mid H)\\
 &=D_0+\mathcal H(Z\mid H)-\mathcal H(d_0).
 \end{split}
\end{equation}
Conditional on $\mathcal A,H$, the variable $Z$ is obtained by adding
the independent $d_0$.  Data processing in
\eqref{eq:conditional-tape-information} therefore yields
$I(Z;\mathcal B\mid H)\le D_0$.

The direct bound $\mathcal H(Z\mid H)=O(\log n)$ would be too
large after summing over scales.  Instead, add the negative tape's
displacement to $Z$.  Their sum is the total common-block gap, whose
entropy can be compared with that of the buffer.  Conditioning on
the negative tape costs at most $D_0$, as just proved:

\begin{align}
 \mathcal H(T,Z+X_B)
 &\ge \mathcal H(T,Z+X_B\mid\mathcal B,H)
  =\mathcal H(T,Z\mid\mathcal B,H)\notag\\
 &\ge \mathcal H(Z\mid H)-D_0
       +\mathcal H(T\mid Z,\mathcal A,\mathcal B,H)\notag\\
 &=\mathcal H(Z\mid H)-D_0+\mathcal H(U\mid d_0).
 \label{eq:buffer-entropy-chain}
\end{align}
For the last equality, the two lists in each chunk determine its common
count.  Knowing these lists and $Z$ therefore fixes both $T-U$ and
$d_0=Z-X_A$.  The buffer pair is independent of the chunks.

The vector $Z+X_B$ is the sum of the first $T$ common gap increments.
The count $T$, like $U$, is independent of the underlying iid tape;
we are not conditioning on the observed chunks in the following
identities.  Consequently
\[
 \mathcal H(T,Z+X_B)=\mathcal H(T)+\mathbf E h_T,
 \qquad
 \mathcal H(U,d_0)=\log n+\mathbf E h_U.
\]
Here $n\le U<2n$ and $10n\le T\le14n-4$, so $T$ takes at most $4n$
values.  Monotonicity of $h_k$ gives
\begin{equation}\label{eq:random-count-entropy-difference}
 \mathcal H(T,Z+X_B)-\mathcal H(U,d_0)
 \le\log4+h_{14n}-h_n.
\end{equation}
Subtract $\mathcal H(d_0)$ in \eqref{eq:buffer-entropy-chain} and use
$\mathcal H(U,d_0)=\mathcal H(d_0)+\mathcal H(U\mid d_0)$.
Together with \eqref{eq:endpoint-information-cost} and
\eqref{eq:random-count-entropy-difference}, this gives explicitly
\[
 D(P_n\Vert Q_n)\le 2D_0+\log4+h_{14n}-h_n.
\]
The bound for $D_0$ is uniform in $n$, so its doubled contribution
is absorbed into $C$.  This proves
Lemma~\ref{lem:entropy-comparison}.
\end{proof}

\subsection{Rare contacts under the reference law}

Comparing opposed paths in one environment and classifying intersection
sites by quenched escape probabilities already appears in the planar
argument of Zerner~\cite[Section~2]{Zerner2007}. We prove the finite-barrier
estimate needed here.

\begin{lemma}\label{lem:reference-contact}
For the reference experiment just defined,
\[
 Q_n(E_{\mathrm{mid}})
 \le\delta_n:=\min\{1,Cn(F_+(n)+F_-(n))\}.
\]
\end{lemma}
\begin{proof}
Fix $(H,Z)$ in its retained support; in particular $H_1,H_3\ge n$.
Write $N=N_*$.  By list-reversal invariance in
Lemma~\ref{lem:exact-bridge}, the placed chronological chunks have
their corresponding raw bridge laws.

Retain the complete positive lower and middle chunks, concatenated
into a path $\alpha$ from height $0$ to its first hit of $N-H_1$.
Retain the negative upper and middle chunks only up to their first
arrival at a departure site of $\alpha$.  A middle contact ensures
that this arrival occurs before the negative middle endpoint at
height $H_3-1$.  The contact site $y$ must therefore satisfy
\begin{equation}\label{eq:reference-contact-height}
 H_3\le y_1\le N-H_1-1.
\end{equation}
Thus each path has made at least $n$ progress in its own signed height
before reaching the contact, and neither has crossed the global
barriers $-1,N$.

We explain the comparison with two paths in one environment before
using any quenched exit probabilities.  The unused positive upper
chunk and negative lower chunk integrate to one.  The remaining four
bridge normalizations are bounded by $p_+^{-2}p_-^{-2}$.  Consecutive
bridges within a sign have disjoint departure layers, so their raw
weights multiply to the raw weight of their concatenation.  Summing
over the negative completion after its retained prefix, and dropping
its remaining bridge requirements, costs at most that prefix's raw
weight: in each fixed environment the total conditional completion
probability is at most one, and this inequality can then be averaged.
This argument allows the completion to revisit its prefix's departure
rows and does not factor their weights independently.
The first-hit chunk boundaries determine the splits, so no
extra multiplicity is introduced.

The retained negative prefix has no departure site in
$\operatorname{Dep}(\alpha)$: its terminal point is the first arrival
there.  Repeated sites within either path cause no difficulty.
Lemma~\ref{lem:departure-transfer} therefore identifies the product
of their raw weights with the weight of the two paths sampled
independently conditional on a common iid environment.  For each
$\alpha$ the first contact fixes the negative prefix uniquely, and
$\alpha$ itself is stopped at a fixed first-hit height.  We may thus
sum these joint prefix weights without overcounting, and extend the
paths as raw trajectories in the common environment.

Only now, in that shared environment, define for interior sites
\[
 Q_y^\omega=P_{y,\omega}(T_{-1}<T_N),
\]
where the barriers are absolute first-coordinate heights.  If
$Q_y^\omega\ge1/2$ at the contact, the positive raw path has visited
such a site after hitting $n$ and before exiting the global slab.
Conditional on the environment, stop it at its first such visit.
The probability of a subsequent exit at $-1$ is at least $1/2$.
Consequently the annealed probability of this visit event is at most
\[
 2P_0(T_n<T_{-1}<T_N)=2(u_n^+-u_N^+).
\]
If $Q_y^\omega<1/2$, the negative marginal gives the analogous bound
$2(u_n^--u_N^-)$.  Indeed the almost-sure finite-height-interval
exit property established in Section~\ref{sec:arrangement} holds
quenched for all these starting sites and barriers outside one
null set.  The probability of the other exit is therefore greater
than $1/2$, with no common lower bound on the transition rows.
Measured from the negative start $N-1$, the
global barriers $N$ and $-1$ have signed relative heights $-1$ and
$N$, respectively.  Lateral translation by $Z$ does not change its
annealed law.

Combining these alternatives with the bounded normalizations gives,
uniformly in $(H,Z)$,
\[
 Q_n(E_{\mathrm{mid}}\mid H,Z)
 \le C\sum_{\sigma\in\{+,-\}}(u_n^\sigma-u_N^\sigma)
 \le CN\bigl(F_+(n)+F_-(n)\bigr),
\]
where the last step is Lemma~\ref{lem:renewal-tail}.  Finally
$\mathbf E N\le12n\mathbf E K$.  Average over the retained joint
marginal of $(H,Z)$ and take the minimum with $1$.
\end{proof}

\subsection{Summing over scales}

\begin{proof}[Proof of Proposition~\ref{prop:entropy-upper}]
For $l\ge2$, take $n=2^{l-2}$ and abbreviate $\delta_l=\delta_n$.
The event $E_{\mathrm{mid}}$ contains the original contact event in
common blocks $(2^l,2^{l+1}]$.  If $\delta_l=0$, absolute continuity
from Lemma~\ref{lem:entropy-comparison} makes its actual probability
zero.  Otherwise the binary relative-entropy inequality gives
\begin{equation}\label{eq:binary-contact-entropy}
 P_n(E_{\mathrm{mid}})\log(1/\delta_l)
 \le C+h_{14n}-h_n.
\end{equation}
For completeness, applying data processing to an event with actual
and reference probabilities $a,b$ gives binary divergence
$a\log(a/b)+(1-a)\log((1-a)/(1-b))$.  It is at least
$a\log(1/b)-\log2$.  Use $b\le\delta_l$ and absorb $\log2$ into $C$.

Width integrability gives
\[
 \sum_{l\ge2}\delta_l<\infty.
\]
Among $2\le l\le J$, at most $C\sqrt J$ indices have
$\delta_l>J^{-1/2}$.  They contribute at most $C\sqrt J$ to $m(J)$.
For every other index with positive $\delta_l$, the logarithm in
\eqref{eq:binary-contact-entropy} is at least $\tfrac12\log J$.
Since $14n\le16n$, monotonicity and bounded overlap give
\[
 \sum_{l=2}^J\bigl(h_{14\cdot2^{l-2}}-h_{2^{l-2}}\bigr)
 \le\sum_{l=2}^J\bigl(h_{2^{l+2}}-h_{2^{l-2}}\bigr)
 \le4h_{2^{J+2}}=O(J),
\]
using \eqref{eq:gap-entropy-growth}.  Summing
\eqref{eq:binary-contact-entropy}, and adding the initial scale, proves
\[
 m(J)\le1+C\sqrt J+\frac{CJ}{\log J}
       =O(J/\log J).
\]
\end{proof}

The estimate uses the same departure contacts as the original
arrangement.  The next section bounds the chance of a first contact
far above a true cut; that bound will force more contact scales than
Proposition~\ref{prop:entropy-upper} permits.
